\writeSectionFrame{\presentationTitle}{Fazit}
\begin{frame}{Anwendungsfälle}
	\begin{itemize}
 		\item Rückgängig machen
 		\item Wiederherstellungspunkte
 	\end{itemize}
\end{frame}

\begin{frame}{Optimierungen}
	\begin{itemize}
 		\item Nur einen Teil des Zustandes speichern
 		\item Inkrementelles Speichern
 	\end{itemize}
\end{frame}

\realworld{Text-Editoren / IDEs}{Zurücksetzen und Wiederherstellen von Dokumenten oder Code}
\note{
    Viele Texteditoren und integrierte Entwicklungsumgebungen (IDEs) nutzen das Memento-Pattern, um Änderungen an Dokumenten oder Code zurückzusetzen oder wiederherzustellen. Jede Bearbeitung oder Änderung wird als Memento gespeichert, sodass der Benutzer zu einem früheren Zustand zurückkehren kann.
}

\realworld{Grafik- und Designsoftware}{Historie von Bearbeitungen und Änderungen}
\note{
    Grafiksoftware wie Adobe Photoshop oder Illustrator verwendet ähnliche Konzepte, um die Historie von Bearbeitungen und Änderungen zu verwalten. Benutzer können zu vorherigen Versionen von Designs oder Bildern zurückkehren.
}

\realworld{Spiele}{Speichern des Spielfortschritts}
\note{
    In vielen Spielen wird das Memento-Pattern verwendet, um den Zustand des Spiels (z.B. Spielfortschritt, Level, Punkte) zu speichern. Dies ermöglicht es den Spielern, ihren Fortschritt zu speichern und später fortzusetzen oder auf frühere Zustände zurückzukehren.
}

\realworld{Transaktionen und Persistenz}{Speichern des Zustandes von Daten vor und nach einer Transaktion}
\note{
    Bei der Implementierung von Transaktionssystemen kann das Memento-Pattern verwendet werden, um den Zustand von Daten vor und nach einer Transaktion zu speichern. Dies hilft, die Datenbank in einen konsistenten Zustand zurückzuversetzen, wenn eine Transaktion fehlschlägt.
}

\vorteil{Kapselung des Zustands}
\note{
    Der Zustand eines Objekts kann gespeichert und wiederhergestellt werden, ohne dass die interne Struktur des Objekts offengelegt werden muss. Dies bewahrt die Kapselung und Integrität des Objekts.
}

\vorteil{Ermöglicht Rückgängig- und Wiederherstellungsoperationen}
\note{
    Es ermöglicht das Implementieren von Undo- und Redo-Funktionalitäten, was besonders in Anwendungen wie Texteditoren und Grafiksoftware nützlich ist.
}

\vorteil{Verlaufsverwaltung}
\note{
    Es erleichtert das Verwalten von Zustandsverläufen und die Rückkehr zu früheren Zuständen, was für Anwendungen wie Spiele und Formulardatenmanagement von Vorteil ist.
}

\vorteil{Trennung von Verantwortlichkeiten}
\note{
    Die Implementierung des Memento-Patterns kann relativ einfach sein und bietet eine klare Struktur zur Verwaltung von Zuständen.
}

\nachteil{Erhöhter Speicherbedarf}
\note{
    Die Speicherung vieler Mementos kann zu erhöhtem Speicherverbrauch führen, insbesondere wenn der Zustand eines Objekts groß oder komplex ist.
}

\nachteil{Komplexität der Verwaltung}
\note{
    Die Verwaltung der gespeicherten Mementos (z.B. bei Undo/Redo) kann komplex werden, besonders wenn eine große Anzahl von Zuständen gespeichert wird.
}

\nachteil{Kopieraufwand}
\note{
    Das Erstellen von Mementos kann zusätzliche Kosten verursachen, insbesondere wenn das Objekt viele Daten enthält oder komplexe Strukturen hat, die tief kopiert werden müssen.
}

\nachteil{Wartungsaufwand}
\note{
    Wenn sich die interne Struktur der Klasse ändert, müssen auch die Mementos und die Klasse, die sie verwaltet, möglicherweise aktualisiert werden. Dies kann zusätzlichen Wartungsaufwand verursachen.
}